\section{广义函数与缓增广义函数}

在本节中，n 表示一个自然数。我们用 \(\mu\) 表示 \(\mathbf{R}^{n}\) 上的勒贝格测度，以及它在 \(\mathbf{R}^{n}\) 的每个局部紧子集上的限制。

我们用

\[
x \cdot y = \sum_ {i = 1} ^ {n} x _ {i} y _ {i}
\]

表示欧几里得空间 \(\mathbf{R}^{n}\) 上的数量积；欧几里得范数记作 \(x \mapsto \|x\|\) （TG, VI, p. 7）。我们回顾，群 \(\mathbf{R}^{n}\) 通过映射 \((x, y) \mapsto \exp(2i\pi x \cdot y)\) 与自身对偶，且勒贝格测度的对偶测度这时与勒贝格测度等同（II, p. 236 的推论 3）。我们用 F （相应地， \(\overline{F}\) ）表示 \(\mathbf{R}^{n}\) 上的傅里叶变换（相应地，傅里叶逆变换）（参见 II, p. 237 的 no. 9）。

对每个 \(\alpha = (\alpha_{i})_{1\leqslant i\leqslant n}\in \mathbf{N}^{n}\) ，我们用 \(X^{\alpha}\) 表示从 \(\mathbf{R}^n\) 到 \(\mathbf{R}\) 中的、由 \(x = (x_{i})_{1\leqslant i\leqslant n}\mapsto x^{\alpha} = \prod_{i = 1}^{n}x_{i}^{\alpha_{i}}\) 所定义的函数。

设 \(\alpha\) 与 \(\beta\) 是 \(\mathbf{N}^n\) 的元素。我们记

\[
| \alpha | = \sum_ {i = 1} ^ {n} \alpha_ {i}, \qquad \binom{\alpha}{\beta} = \prod_ {i = 1} ^ {n} \binom{\alpha_ {i}}{\beta_ {i}}.
\]

引理 1. --- 设 U 是 \(\mathbf{R}^{n}\) 的开子集。

\begin{enumerate}
\def\labelenumi{\alph{enumi})}
\item
  设 K 是 U 的紧子集，V 是 K 在 U 中的开邻域。存在 \(\varphi \in \mathcal{C}^{\infty}(\mathrm{U})\) ，其紧支撑含于 V，使得 \(0 \leqslant \varphi \leqslant 1\) ，且 \(\varphi(x) = 1\) 对所有 \(x \in \mathrm{K}\) 成立；
\item
  对 U 的每个局部有限开覆盖，存在一个由 \(C^{\infty}\) 类函数组成、从属于它的单位分解。
\end{enumerate}

这由 VAR, R1, p. 40, 5.3.6 得出。

命题 1. --- 设 U 是 \(\mathbf{R}^{n}\) 的开子集，\(k\in\mathbf{N}\) 。设 E、F、G 是拓扑向量空间，b: E×F→G 是连续双线性映射。设 f 与 g 分别是 U 到 E 与 F 中的 \(\mathbf{C}^{k}\) 类函数。那么映射 h: x↦b(f(x),g(x)) 在 U 上是 \(\mathbf{C}^{k}\) 类的；此外，对每个 \(\alpha\in\mathbf{N}^{n}\) （满足 \(|\alpha|\leqslant k\) ）以及

对每个 \(x \in \mathrm{U}\) ，我们有

\[
\partial^{\alpha}h(x) = \sum_{\substack{\beta \in \mathbf{N}^{n}\\ \beta \leqslant \alpha}}\binom {\alpha}{\beta}b\big(\partial^{\beta}f(x),\partial^{\alpha -\beta}g(x)\big).
\]

映射 h 是映射 \((f,g)\) : \(U \rightarrow E \times F\) （它是 \(C^{k}\) 类的）与映射 b: \(E \times F \rightarrow G\) （它是 \(C^{\infty}\) 类的）的复合。因而它是 \(C^{k}\) 类的。其偏导数的表达式由莱布尼茨公式得出（FVR, I, p. 28，命题 2，参见 A, III, p. 122，推论）。

我们回顾，分离的局部凸拓扑向量空间若为桶型空间且每个有界子集都相对紧，则称为蒙泰尔空间（EVT, IV, p. 18，定义 4）。

我们还回顾，若 E 与 F 是局部凸空间，且 E 是有界型的（EVT, III, p. 12，定义 1），则线性映射 \(u\colon\mathrm{E}\to\mathrm{F}\) 连续当且仅当 E 的每个有界子集在 \(u\) 下的像在 F 中有界（EVT, III, p. 11，命题 1, (iiibis)；当 F 是可赋范时，一般情形可形式地由此推出，参见 EVT, II, p. 7，命题 5, b)）。

